\section{Conclusion}
\label{sec:conclusion}

Typed decision models bring classification, scoring and Boolean judgments into one interface whose questions are specified at call time. For new decisions, they address the goal that routing surveys state: a decision component that adapts without retraining. The routing approaches reviewed here still need candidate-specific evidence for new models, and risk targets still need labeled evidence. Independent studies report valid outputs and large savings against LLM judges, although a cheap generative cascade can be cheaper inside an agent loop. They also show accuracy gaps on hard tasks, calibration that varies by workload, and decisions that added context can move. We recommend three steps: 1) compare a typed model with the alternatives in Table~\ref{tab:starting}, including a label-probability readout on an open model and a cheap generative cascade; 2) measure both under the acceptance rule the deployment will use, on labeled data from its own traffic; and 3) record which labels, profiles and warmup batches can be reused when models, questions or constraints change.

\noindent\textbf{Limitations.} Most work specific to Jev is recent and not peer reviewed, and results can change between versions of a paper or of the model. Jev's weights and training data are closed, so black-box measurements cannot separate model behavior from provider-side processing. The JevBench comparisons use one results snapshot without reported uncertainty estimates. Our search covered public English-language sources and may have missed evaluations that have not appeared as papers or documented repositories. \rev{Work posted after September 30, 2026, is not covered systematically.}
